\documentclass[12pt]{article}
\usepackage[T1]{fontenc}
\usepackage[utf8]{inputenc}
\usepackage{lmodern}
\usepackage{textcomp}
\usepackage{enumitem}
\usepackage{geometry}
\geometry{margin=1in}
\begin{document}
\begin{center}
Answer Key - Political Science - Graduate Level Assessment 24 \
\small (Answers are ordered exactly as in the question paper.)
\end{center}

\section*{Multiple Choice}
\begin{enumerate}[label=\arabic*.]
\item \textbf{Answer:} A
\\ \textit{Options:} A) US banks funded the allies much more than Germany and Austria-Hungary; B) US banks funded Germany and Austria-Hungary much more than the allies; C) Wilson secretly favoured the German position; D) None of the above
\\ \textit{Note:} The correct answer highlights that the substantial financial support provided by US banks to the Allies created a conflict of interest that ultimately undermined America's purported neutrality, as the economic ties made it increasingly difficult for the US to remain impartial in the face of escalating hostilities.
\item \textbf{Answer:} A
\\ \textit{Options:} A) Voluntary reliance on an external power for security; B) Willful openness to colonization; C) Cultural imperialism; D) Open advocacy of imperialism for economic gain
\\ \textit{Note:} The phrase 'empire by invitation' refers to a situation where a state voluntarily seeks the assistance of a more powerful external entity for security purposes, thereby establishing a relationship that resembles an empire without coercion.
\item \textbf{Answer:} B
\\ \textit{Options:} A) Armed conflict; B) Anarchy; C) Institutions; D) No common language
\\ \textit{Note:} The correct answer is "anarchy" because realists argue that the international system lacks a central authority to enforce rules and maintain order, distinguishing it fundamentally from the domestic system, where a government provides structure and governance.
\item \textbf{Answer:} D
\\ \textit{Options:} A) James Madison; B) Abraham Lincoln; C) Woodrow Wilson; D) Thomas Jefferson
\\ \textit{Note:} The phrase "Peace, commerce, and honest friendship with all nations, entangling alliances with none" reflects Thomas Jefferson's foreign policy philosophy articulated in his inaugural address, emphasizing non-interventionism and the avoidance of permanent alliances.
\item \textbf{Answer:} B
\\ \textit{Options:} A) the Senate.; B) the president.; C) the Secretary of State.; D) the chairman of the Joint Chiefs.
\\ \textit{Note:} The president has the constitutional authority to recognize foreign governments, a power that is essential for conducting foreign relations and shaping U.S. diplomacy.
\end{enumerate}

\section*{True / False}
\begin{enumerate}[label=\arabic*.]
\setcounter{enumi}{5}
\item \textbf{Answer:} False
\\ \textit{Options:} A) True; B) False
\\ \textit{Note:} The Vietnam War illustrated that factors such as guerrilla tactics, local support, and political will can undermine technological superiority, leading to military failure despite advanced weaponry.
\item \textbf{Answer:} True
\\ \textit{Options:} A) True; B) False
\\ \textit{Note:} The phrase 'indiscriminate attack' is associated with a violation of international humanitarian law and ethical principles, undermining the legitimacy of military operations by failing to distinguish between combatants and civilians, which is essential for maintaining moral and legal justification in warfare.
\item \textbf{Answer:} False
\\ \textit{Options:} A) True; B) False
\\ \textit{Note:} The statement is false because the United States actively engaged in the Third World during the Cold War by providing economic and military assistance, supporting anti-communist regimes, and intervening in various conflicts to counter Soviet influence.
\item \textbf{Answer:} True
\\ \textit{Options:} A) True; B) False
\\ \textit{Note:} Systemic theories emphasize that the structure and distribution of power among states in the international system fundamentally shape their foreign policy decisions, as they must navigate the constraints and opportunities presented by their relative power positions.
\item \textbf{Answer:} True
\\ \textit{Options:} A) True; B) False
\\ \textit{Note:} The cyber-security discourse gained traction in the 1970 s as the rise of personal computing and the internet prompted policymakers and scholars to address the emerging risks associated with digital information and infrastructure.
\end{enumerate}

\section*{Long Form Response}
\begin{enumerate}[label=\arabic*.]
\setcounter{enumi}{10}
\item \textbf{Answer:} The United Nations Security Council (UNSC) is structured with five permanent members- China, France, Russia, the United Kingdom, and the United States-each possessing veto power, alongside ten rotating members elected for two-year terms without such authority. This arrangement reflects a balance between the interests of major global powers and the representation of a broader international community. The veto power held by permanent members allows them to block substantive resolutions, which can lead to significant implications for global governance, often resulting in inaction during crises when consensus is not achieved. Conversely, the rotating members provide a platform for diverse perspectives but lack the influence to challenge the decisions of the permanent members, highlighting a structural imbalance that can affect the effectiveness and legitimacy of the UNSC in addressing international security issues. Thus, while the structure aims to maintain stability, it can also hinder the Council's responsiveness to global challenges.
\\ \textit{Note:} The answer accurately identifies the structural composition of the UNSC, emphasizing the significant influence of the five permanent members through their veto power, while also acknowledging the role of rotating members, which collectively illustrates the tension between major power interests and broader international representation in global governance.
\item \textbf{Answer:} New critical schools of security studies have significantly reshaped the field by challenging traditional notions of security that prioritize state-centric and military- oriented perspectives. These schools advocate for a broader understanding of security that includes human security, environmental concerns, and the role of non-state actors, thereby redefining the boundaries of the discipline. Their relationship with international relations is complex, as they both draw from and critique conventional IR theories, pushing the discipline toward inclusivity and reflexivity. This evolving dynamic suggests that security studies occupies a unique position between academic inquiry and practical application, often bridging the gap between theoretical frameworks and real-world challenges faced by practitioners. Ultimately, the emergence of these new critical perspectives has profound implications for the identity of security studies, positioning it as a more interdisciplinary and socially engaged field within political science.
\\ \textit{Note:} correct because it emphasizes how new critical schools of security studies expand the traditional focus of the discipline by incorporating diverse perspectives on security that address human, environmental, and non-state actor concerns, which in turn complicates and enriches the relationship with international relations and the overall identity of the field.
\end{enumerate}

\vfill
\noindent\textit{End of Answer Key}
\end{document}